%     AMS-LaTeX v.2 template for use with amsart
%
%     Remove any commented or uncommented macros you do not use.

\documentclass[%
  aps,%
  prc,%
  showpacs,%
  superscriptaddress,%
  onecolumn,%
  notitlepage,%
  12pt,%
  floatfix,%
  % endfloats,%
  amsmath,%
  amssymb,%
  %preprintnumbers,%
  %linenumbers,%
]{revtex4-1}

\usepackage{amscd}
\usepackage{physics}

\newcommand{\clebsch}[6]{(#1 #2 #3 #4 \vert #5 #6)}
\newcommand{\threej}[6]{
  \renewcommand\arraystretch{0.6667}\begin{pmatrix} #1&#3&#5\\#2&#4&#6 \end{pmatrix}\renewcommand\arraystretch{1}
}
% \newcommand{\threej}{\begin{pmatrix} j'&L&j\\-m'&M&m \end{pmatrix}}
\newcommand{\me}{\langle \beta j' m' \rvert T_{LM} \lvert \alpha j m \rangle}
\newcommand{\rme}{\langle \beta j' \rvert\vert T_L \vert\lvert \alpha j \rangle}
\newcommand{\jphat}{\hat\jmath'}
\newcommand{\hide}{\hphantom}

\begin{document}

\title{Calculating M-scheme matrix elements from coupled two-body RMEs}

\author{Patrick J. Fasano}
\affiliation{Department of Physics, University of Notre Dame, Notre Dame, Indiana 46556-5670, USA}

\author{Pieter Maris}
\affiliation{Department of Physics and Astronomy, Iowa State University, Ames, Iowa 50011-3160, USA}

\date{\today}

\maketitle


Numerous conventions in the literature exist for the definition of the reduced
matrix element in the Wigner-Eckart theorem. Here we derive the relationship
between matrix elements (of a spherical tensor operator) in the uncoupled basis
to reduced matrix elements in the coupled basis. We adopt the following convention
for the Clebsch-Gordan coefficients and Wigner 3J symbols \cite{suhonen2007:nucleons-nucleus}
\begin{equation}
  (j_a m_a j_b m_b \vert J M) = (-1)^{j_b-j_a-M}\hat{J} \threej{j_a}{m_a}{j_b}{m_b}{J}{-M}
\end{equation}
where \(\hat{J} = \sqrt{2J+1}\).

First, we relate the coupled and uncoupled antisymmetrized two-particle bases:
\begin{align}
  \ket{a m_a b m_b}_\mathrm{AS} &= \sum_J (j_a m_a j_b m_b \vert J M) \ket{a b; J M}_\mathrm{AS}\\
  &=\sum_J (-1)^{j_b-j_a-M}\hat{J} \threej{j_a}{m_a}{j_b}{m_b}{J}{-M} \ket{a b; J M}_\mathrm{AS}
\end{align}

Next we insert these into the matrix element:
\begin{align*}
  & \bra{a m_a b m_b}\mathcal{O}\ket{c m_c d m_d} \\
    &= \qty(\sum_{J'} (-1)^{j_b-j_a-M'}\hat{J'} \threej{j_a}{m_a}{j_b}{m_b}{J'}{-M'} \bra{a b; J' M'}_\mathrm{AS})
      \mathcal{O} \\
    &\qquad \times \qty(\sum_J (-1)^{j_d-j_c-M}\hat{J} \threej{j_c}{m_c}{j_d}{m_d}{J}{-M} \ket{c d; J M}_\mathrm{AS})\\
    &= (-1)^{j_b -j_a} (-1)^{j_d-j_c} (-1)^{-M'-M} \\
    &\qquad \times \sum_{J',J} \hat{J'} \hat{J} \threej{j_a}{m_a}{j_b}{m_b}{J'}{-M'}
      \threej{j_c}{m_c}{j_d}{m_d}{J}{-M} \bra{ab;J'M'}\mathcal{O}\ket{cd;JM}_\mathrm{AS}
\end{align*}

We use the Wigner-Eckart theorem to write the matrix element in terms of reduced matrix elements
\begin{equation*}
  \bra{\alpha J' M'}\mathcal{O}_{J_0 M_0}\ket{\beta J M}
    = (-1)^{J'-M'}\threej{J'}{-M'}{J_0}{M_0}{J}{M} \langle{\alpha J'}\vert\vert\mathcal{O}_{J_0}\vert\vert{\beta J}\rangle
\end{equation*}


Finally, we express the uncoupled matrix element in terms of the coupled reduced matrix element:
\begin{align*}
  & \bra{a m_a b m_b}\mathcal{O}_{J_0 M_0}\ket{c m_c d m_d} \\
  &= (-1)^{j_b -j_a} (-1)^{j_d-j_c} (-1)^{-M'-M} \\
    &\qquad \times \sum_{J',J} \hat{J'} \hat{J} \threej{j_a}{m_a}{j_b}{m_b}{J'}{-M'}
      \threej{j_c}{m_c}{j_d}{m_d}{J}{-M} \bra{ab;J'M'}\mathcal{O}\ket{cd;JM}_\mathrm{AS}\\
  &= (-1)^{j_b -j_a} (-1)^{j_d-j_c} (-1)^{-M'-M} \\
    &\qquad \times \sum_{J',J} (-1)^{J'-M'} \hat{J'} \hat{J} \threej{j_a}{m_a}{j_b}{m_b}{J'}{-M'}
      \threej{j_c}{m_c}{j_d}{m_d}{J}{-M} \threej{J'}{-M'}{J_0}{M_0}{J}{M} \langle{ab; J'}\vert\vert\mathcal{O}_{J_0}\vert\vert{cd; J}\rangle\\
\end{align*}

If we assume that $j_a,j_b,j_c,j_d \in \mathbb{Z}+\frac{1}{2}$ and that
$J',M',J,M,J_0,M_0 \in \mathbb{Z}$, we can simplify and rearrange the phase
factors to a form which can be implemented in \texttt{mfdn-transitions}:

\begin{multline}
  \bra{a m_a b m_b}\mathcal{O}_{J_0 M_0}\ket{c m_c d m_d} = \\
  (-1)^{j_b-j_a + j_d-j_c - M}
  \sum_{J'} (-1)^{J'} \hat{J'} \threej{j_a}{m_a}{j_b}{m_b}{J'}{-M'}
  \sum_{J} \hat{J} \threej{j_c}{m_c}{j_d}{m_d}{J}{-M} \threej{J'}{-M'}{J_0}{M_0}{J}{M} \langle{ab; J'}\vert\vert\mathcal{O}_{J_0}\vert\vert{cd; J}\rangle
\end{multline}
% \bibliographystyle{apsrevm}
\bibliography{master,mc,theory,expt,data,books,proc,misc}

\end{document}

%-----------------------------------------------------------------------
% End of amsart-template.tex
%-----------------------------------------------------------------------
